\section{Derivation of Mode Matching TM00 Beams}
%%
%A Gaussian beam can be described by the complex beam parameter \cite{Bond2016}, $q$,
%\begin{equation}
%q(z) = z - z_0 + iz_R,
%\end{equation}
%where $z$ is the distance along the optical axis, $z_0$ is the beam waist location, and $z_R$ is the Rayleigh length of the beam. This is a useful way of describing a Gaussian beam because the defocus and width of the beam at any point along $z$ can be calculated from it. Taking the reciprocal of $q(z)$ gives
%\begin{equation}
%\frac{1}{q(z)} = S(z) -i\frac{\lambda}{\pi\omega^2(z)},
%\label{eq:1/q}
%\end{equation}   
%where $S(z)$ is the defocus at $z$, $\lambda$ is the lasers wavelength and $w(z)$ is the width of the beam at $z$. From equation \ref{eq:1/q} we find
%\begin{align}
%S(z) &= \Re\left( \frac{1}{q} \right),\\
%w(z) & =  \sqrt{\frac{\lambda}{-\pi\Im\left(\frac{1}{q}\right)}}.
%\end{align}
%To calculate how the beam parameter $q_0$ is effected by propagation $\Delta z$ and reflection from a curved mirror with power $S$, the coefficient matrices $\mathbf{P}$ and $\mathbf{R}$ can be applied to give the resultant beam parameter $q_1$
%\begin{align}
%\mathbf{M} = \begin{pmatrix} A &B\\	C &D \end{pmatrix},\quad \mathbf{P} &= \begin{pmatrix} 1 &\Delta z\\	0 &1 \end{pmatrix},\quad \mathbf{R} = \begin{pmatrix} 1 & 0\\	-S & 1 \end{pmatrix},\\
%q_1 &= \frac{Aq_0 + B}{Cq_0 + D}.
%\end{align}
%\paragraph{}
%If a beams angle on incidence on a curved mirror is not zero, the reflected beam with have astigmatism. The location of the beam waist will be different for the tangential plane (plane containing the optical axis of the mirror and the object) and the sagittal (plane perpendicular to the tangential plane). The shape of the mirror can be described by the tangential radius of curvature, $R_t$, and the sagittal radius of curvature, $R_s$. If the beam has an angle of incidence $\theta$ on a mirror with a radius of curvature $R$, then the radius of curvature experienced by the tangential and sagittal beams is \cite{MonkLight,Bond2016}
%\begin{align}
%R_t &= R\cos\theta \\
%R_s &= \frac{R}{\cos\theta}
%\end{align}
%so the optical power is changed by
%\begin{align}
%S_t &= \frac{2}{R_t} = \frac{S}{\cos\theta}\\
%S_s &= \frac{2}{R_s} = S\cos\theta.
%\end{align}
%The radius of curvature of a mirror is related to the optical power of the mirror by $S_m= \frac{2}{R_m}$, and the defocus of a wavefront is related to the curvature of the wavefront by $R_w = \frac{1}{R_w}$.
%\paragraph{}
%The electric field, $E$, of a TM00 Gaussian beam can be written as
%\begin{equation}
%E(x,y,z) = E_x(x,z)E_y(y,z)
%\end{equation}
%where 
%\begin{align}
%E_x(x,z) = \frac{1}{\sqrt{q_x(z)}}\exp\left(-ik\frac{x^2}{2q_x(z)}\right),\\
%E_y(y,z) = \frac{1}{\sqrt{q_y(z)}}\exp\left(-ik\frac{y^2}{2q_y(z)}\right).
%\end{align}
%For a beam without astigmatism, $x^2 + y^2 = r^2$ and $q_x = q_y = q$ so:
%\begin{equation}
%E(r,z) = \frac{1}{q(z)}\exp\left(-ik\frac{r^2}{2q(z)}\right).
%\end{equation}
%%]
The \textcolor{red}{check this naming} electric field, $E$, of a TM00 Gaussian beam can be written as
\begin{equation}
E(x,y,z) = E_x(x,z)E_y(y,z)
\end{equation}
where 
\begin{align}
E_x(x,z) = \frac{1}{\sqrt{q_x(z)}}\exp\left(-ik\frac{x^2}{2q_x(z)}\right),\\
E_y(y,z) = \frac{1}{\sqrt{q_y(z)}}\exp\left(-ik\frac{y^2}{2q_y(z)}\right).
\end{align}
For a beam without astigmatism, $x^2 + y^2 = r^2$ and $q_x = q_y = q$ so:
\begin{equation}
E(r,z) = \frac{1}{q(z)}\exp\left(-ik\frac{r^2}{2q(z)}\right).
\end{equation}
The mode matching, $\mathcal{M}$, of two beams can be expressed by overlap of the two fields
\begin{equation}
\mathcal{M} = \frac{\left|\int_{A} E_1^*E_2 dA \right|}{\int_{A} \left|E_1 \right|^2 dA \int_{A} \left|E_2 \right|^2 dA}.
\label{eq:mode_matching_general}
\end{equation}
The numerator of equation \ref{eq:mode_matching_general} for TM00 beams can be written as:
\begin{align}
\left|\int_{A} E_1^*E_2 dA \right| &= \int_{A} E_1^*E_2 dA\int_{A} E_1E_2^* dA \\
&= \int^{-\infty}_{-\infty} E_{x1}^*E_{x2} dx \int^{-\infty}_{-\infty} E_{y1}^*E_{y2} dy \int^{-\infty}_{-\infty} E_{x1}E_{x2}^* dx \int^{-\infty}_{-\infty} E_{y1}E_{y2}^* dy.
\label{eq:numerator}
\end{align}
Since the calculation is the same for the $x$ and $y$ coordinates, we only need to show the solution for the $x$. The $x$ terms of equation \ref{eq:numerator} can be written as:
\begin{align}
\int^{-\infty}_{-\infty} E_{x1}^*E_{x2} dx &\int^{-\infty}_{-\infty} E_{x1}E_{x2}^* dx = \int_{-\infty}^{\infty} \frac{1}{\sqrt{q_{x1}^*}}\exp\left(\frac{ikx^2}{2q_{x1}^*}\right)\frac{1}{\sqrt{q_{x2}}}\exp\left(\frac{-ikx^2}{2q_{x2}}\right)dx\\
&\times\int_{-\infty}^{\infty} \frac{1}{\sqrt{q_{x1}}}\exp\left(\frac{-ikx^2}{2q_{x1}}\right)\frac{1}{\sqrt{q_{x2}^*}}\exp\left(\frac{ikx^2}{2q_{x2}^*}\right)dx\\
&=\frac{1}{\left|q_{x1}\right|\left|q_{x2}\right|}\left(\int_{-\infty}^{\infty}\exp\left(\frac{ikx^2}{2q_{x1}^*}\right)\exp\left(\frac{-ikx^2}{2q_{x2}^*}\right)dx \times c.c.\right) \\
&=\frac{1}{\left|q_{x1}\right|\left|q_{x2}\right|}\left( \int_{-\infty}^{\infty} \exp\left(\frac{ikx^2}{2q_{x1}^*} -  \frac{-ikx^2}{2q_{x2}^*}\right) dx \times c.c.\right) \\
&=\frac{1}{\left|q_{x1}\right|\left|q_{x2}\right|}\left( \int_{-\infty}^{\infty} \exp\left(\frac{ikx^2\left(q_{x2} - q{x1}^*\right)}{2q_{x1}^*q_{x2}}\right) dx \times c.c.\right).
\label{eq:numerator_x_coords}
\end{align}
Using the identity
\begin{equation}
\mathcal{I} = \int_{-\infty}^{\infty} \exp\left(iax^2\right)dx = \frac{i\pi}{a}
\end{equation}
and labelling 
\begin{equation}
a = \frac{k\left(q_{x2} - q_{x1}*\right)}{2q_{x1}^*q_{x2}}
\end{equation}
allows equation \ref{eq:numerator_x_coords} to be simplified to 
\begin{align}
\frac{1}{\left|q_{x1}\right|\left|q_{x2}\right|}\left( \int_{-\infty}^{\infty} \exp\left(\frac{ikx^2\left(q_{x2} - q{x1}^*\right)}{2q_{x1}^*q_{x2}}\right) dx \times c.c.\right) &= \frac{1}{\left|q_{x1}\right|\left|q_{x2}\right|}\sqrt{\frac{i\pi\left(2q_{x1}^*q_{x2}\right)}{k\left(q_{x2} - q_{x1}^*\right)}}\sqrt{\frac{-i\pi\left(2q_{x1}q_{x2}^*\right)}{k\left(q_{x2}^* - q_{x1}\right)}}\\
& = \frac{1}{\left|q_{x1}\right|\left|q_{x2}\right|} \sqrt{ \frac{4\pi\left|q_{x1}\right|^2\left|q_{x2}\right|^2}{k^2\left|q_{x2} - q_{x1}^*\right|^2} } \\
&= \frac{2\pi}{k}\frac{1}{\left|q_{x2} - q_{x1}^*\right|}.
\end{align}
This means that the numerator of equation \ref{eq:mode_matching_general} is
\begin{equation}
\frac{4\pi^2}{k^2}\frac{1}{\left|q_{x2} - q_{x1}^*\right|\left|q_{y2} - q_{y1}^*\right|}.
\end{equation}
Since $E_x$ and $E_y$ are orthogonal, $\left|E\right|^2 = \left|E_x\right|^2\left|E_y\right|^2 $, an expression for the denominator of equation \ref{eq:mode_matching_general} can be found by:
\begin{align}
\int_A \left|E\right|^2 dA &= \int_{-\infty}^{\infty}\int_{-\infty}^{\infty} \left|E_x\right|^2\left|E_y\right|^2 dx dy \\
&= \int_{-\infty}^{\infty} \frac{1}{\left|q_x\right|}\exp\left(\frac{ikx^2\left(q_x - q_x^*\right)}{2q_xq_x^*}\right) \times \text{y term}.
\end{align}
Using a similar approach as in the calculation of the numerator terms:
\begin{align}
\int_A \left|E\right|^2 dA &= \frac{1}{\left|q_x\right|}\sqrt{\frac{i2\pi\left|q_x\right|^2}{k\left(q_x-q_x^*\right)}}\frac{1}{\left|q_y\right|}\sqrt{\frac{i2\pi\left|q_y\right|^2}{k\left(q_y-q_y^*\right)}} \\
&=-\frac{2\pi}{k}\frac{1}{\sqrt{\left(q_x-q_x^*\right)\left(q_y-q_y^*\right)}}
\end{align}
we get the mode matching $\mathcal{M}$ to be
\begin{align}
\mathcal{M} &= \frac{\frac{4\pi}{k^2}}{-\frac{2\pi}{k}\times-\frac{2\pi}{k}}\frac{\sqrt{\left(q_{x1}-q_{x1}^*\right)\left(q_{y1}-q_{y1}^*\right)}\sqrt{\left(q_{x2}-q_{x2}^*\right)\left(q_{y2}-q_{y2}^*\right)}}{\left|q_{x2} - q_{x1}^*\right|\left|q_{y2} - q_{y1}^*\right|}\\
&=\frac{\sqrt{\left(q_{x1}-q_{x1}^*\right)\left(q_{y1}-q_{y1}^*\right)\left(q_{x2}-q_{x2}^*\right)\left(q_{y2}-q_{y2}^*\right)}}{\left|q_{x2} - q_{x1}^*\right|\left|q_{y2} - q_{y1}^*\right|}.
\end{align}
\section{Effect of Error in the Radius of Curvature of a Mirror on the defocus of a Beam}
The error in the optical power of a mirror due to an error in the radius of curvature of is
\begin{align}
S &= \frac{2}{R_m}\\
dS &= \frac{-2}{R_m^2}dR.
\end{align}
The error in the beams defocus is the same as the error in the radius of curvature of the the mirror. The effect this has on the beam parameter $q$ can be found by defining a contour of modes distance $dq$ from $q$ i.e. $q + \left|dq\right|e^{i\theta}$, where $dq$ is
\begin{align}
q &= \frac{q_i}{Sq_i + 1} \\
dq &= \frac{q_i^2}{(-Sq_i + 1)^2} dS.
\end{align}
This gives you a family of modes which will contain the modes due to an error in the radius of curvature of the mirror, but since $q$ is a complex number the contour you need to define to encapsulate the concept of an error in $q$ also contains modes which are have different beam sizes as well as different defocuses. To do the calculation this way is cumbersome, it is much more straight forward to compute the effect of having a mirror ROC error by evaluating at every mirror
\begin{equation}
\begin{pmatrix} 1 & 0\\	\pm dS &1 \end{pmatrix} \begin{pmatrix} 1 & 0\\	-S &1 \end{pmatrix} \begin{pmatrix} q \\	1 \end{pmatrix} = \begin{pmatrix} q \\(S\pm dS)q + 1 \end{pmatrix} 
\end{equation}
and then divide the top and bottom entries of the $\begin{pmatrix} q \\	(S\pm dS)q + 1 \end{pmatrix} $ vector. 